
\medskip
\phantomsection\label{prop:sga3-VIB-2.5}%
\noindent\textbf{Proposition 2.5.---}%
\footnote{\label{note:sga3-VIB-20}Editor's note: we have expanded
the statement and proof, and in \textup{\SGARefLink{concl:sga3-VIB-2.5-i}{(i)}} we have weakened the
hypothesis on \(H\), replacing ``of finite presentation'' by
``of finite type.''}
\emph{Let \(u:G\to H\) be a morphism of \(S\)-groups.}

\begin{enumerate}
\item[\textup{(i)}]\phantomsection\label{concl:sga3-VIB-2.5-i} \emph{Suppose that \(G\) (respectively, \(H\))
  is of finite presentation and flat (respectively, of finite type)
  over \(S\) at the points of its identity section. Then the sets}
  \[
    U_{\mathrm{plat}}
      \supset U_{\mathrm{lisse}}
      \supset U_{\mathrm{\acute{e}t.}},
  \]
  \emph{formed by the points \(s\in S\) such that \(u_s\) is flat
  (respectively, smooth, or \'etale), are open in \(S\).}

  \emph{If, moreover, \(G\) (respectively, \(H\)) is locally of
  finite presentation and flat (respectively, locally of finite
  type) over \(S\), then the set \(V_{\mathrm{plat}}\)
  (respectively, \(V_{\mathrm{lisse}}\) or
  \(V_{\mathrm{\acute{e}t.}}\)) of points of \(G\) at which \(u\)
  is flat (respectively, smooth, or \'etale) is the inverse image
  under \(\pi_G\) of \(U_{\mathrm{plat}}\) (respectively,
  \(U_{\mathrm{lisse}}\) or \(U_{\mathrm{\acute{e}t.}}\)).}

\item[\textup{(ii)}]\phantomsection\label{concl:sga3-VIB-2.5-ii} \emph{Suppose that, for every \(s\in S\),
  \(G_s\) is locally of finite type over \(\kappa(s)\) (a condition
  satisfied if \(G\) is of finite type over \(S\) at the points of
  the identity section \textup{(\SGARefLink{lem:sga3-VIB-1.3.1.1}{1.3.1.1})}), and that \(u\) is
  locally of finite type (respectively, locally of finite
  presentation) at the points of the identity section of \(G\).
  Then the sets}
  \[
    U_{\mathrm{l.q.f.}}\supset U_{\mathrm{n.r.}},
  \]
  \emph{formed by the \(s\in S\) such that \(u_s\) is locally
  quasi-finite (respectively, unramified), are open in \(S\).}

  \emph{If, moreover, \(u\) is locally of finite type
  (respectively, locally of finite presentation), then the set
  \(V_{\mathrm{l.q.f.}}\) (respectively,
  \(V_{\mathrm{n.r.}}\)) of points of \(G\) at which \(u\) is
  quasi-finite (respectively, unramified) is the inverse image under
  \(\pi_G\) of \(U_{\mathrm{l.q.f.}}\) (respectively,
  \(U_{\mathrm{n.r.}}\)).}
\end{enumerate}

We prove \textup{\SGARefLink{concl:sga3-VIB-2.5-i}{(i)}}. First note (\SGARefLink{lem:sga3-VIB-1.3.1.1}{1.3.1.1}) that, for every
\(s\in S\), \(G_s\) is locally of finite type over \(\kappa(s)\).
Let \(Y\) be the open subset of \(H\) consisting of the points at
which \(\pi_H\) is of finite type, and let \(X_1\) be the open
subset of \(u^{-1}(Y)\) consisting of the points at which
\(\pi_G\) is of finite presentation. By EGA~IV\(_3\), 11.3.1, the
set \(X\) of points of \(X_1\) at which \(\pi_G\) is of finite
presentation and flat is open in \(X_1\), hence in \(G\).

Let
\[
  u_X:X\longrightarrow Y
\]
be the morphism induced by \(u\). Since \(Y\) (respectively, \(X\))
contains the identity section of \(H\) (respectively, \(G\)), the
restriction \(\pi_X\) of \(\pi_G\) to \(X\) is surjective, and the
morphism
\[
  S\longrightarrow X
\]
induced by \(\epsilon_G\) is an \(S\)-section of \(X\).

Consider the sets
\[
  W_{\mathrm{plat}}
    \supset W_{\mathrm{lisse}}
    \supset W_{\mathrm{\acute{e}t.}},
\]
consisting of the points of \(X\) at which \(u_X\) is flat
(respectively, smooth, or \'etale). It is clear that
\(W_{\mathrm{lisse}}\) and \(W_{\mathrm{\acute{e}t.}}\) are open
in \(X\). On the other hand, since \(\pi_Y\)
is locally of finite type and
\[
  \pi_X=\pi_Y\circ u_X
\]
is locally of finite presentation, EGA~IV\(_1\), 1.4.3(v), shows
that \(u_X\) is locally of finite presentation. Consequently, the
flatness locus \(W_X\) of \(u_X\) is open in \(X\)
(EGA~IV\(_3\), 11.3.1).

Let \(x\in X\) and set \(s=\pi_X(x)\). Then, by
EGA~IV\(_2\), 11.3.10,%
\begingroup
\renewcommand{\thefootnote}{\fnsymbol{footnote}}%
\footnote[1]{\label{note:sga3-VIB-2.5-EGA-volume}Translator's note:
the source citation is retained. Section 11 occurs in EGA~IV\(_3\),
not EGA~IV\(_2\); the cited result is EGA~IV\(_3\), 11.3.10.}%
\endgroup{}
combined with EGA~IV\(_4\), 17.5.1 (respectively, 17.6.1),
\(x\) belongs to \(W_{\mathrm{plat}}\) (respectively,
\(W_{\mathrm{lisse}}\) or \(W_{\mathrm{\acute{e}t.}}\)) if and
only if \(u_s\) is flat (respectively, smooth, or \'etale) at the
point \(x\), or, equivalently by \hyperref[prop:sga3-VIB-1.3]{1.3}, if and only if \(u_s\) is
flat (respectively, smooth, or \'etale).

Consequently, for
\[
  \text{``\(\flat=\) flat, smooth, or \'etale,''}
\]
we have
\[
  U_\flat=\epsilon_G^{-1}(W_\flat),
  \qquad
  W_\flat=\pi_X^{-1}(U_\flat).
\]
Since \(W_\flat\) is open in \(X\), hence in \(G\), the first
equality shows that \(U_\flat\) is open in \(S\).

The second assertion of \textup{\SGARefLink{concl:sga3-VIB-2.5-i}{(i)}} follows from what we have just
seen, for then
\[
  Y=H,\qquad X=G,\qquad
  V_\flat=W_\flat=\pi_G^{-1}(U_\flat).
\]

We prove assertion \textup{\SGARefLink{concl:sga3-VIB-2.5-ii}{(ii)}}. Let
\[
  V_{\mathrm{l.t.f.}}\supset V_{\mathrm{l.q.f.}}
\]
be the open subsets of \(G\) consisting of the points at which
\(u\) is of finite type (respectively, quasi-finite). Set
\[
  X=V_{\mathrm{l.t.f.}},
\]
and let \(u_X\) denote the restriction of \(u\) to \(X\). By
hypothesis, \(X\) contains \(\epsilon(S)\). Let \(x\in X\) and set
\(s=\pi_X(x)\).

If \(u\) is quasi-finite at \(x\), then, by base change, \(u_s\) is
quasi-finite at \(x\); and since \(G_s\) is assumed to be locally
of finite type, \(u_s\) is locally quasi-finite by \SGARefLink{cor:sga3-VIB-1.3.1}{1.3.1}.

Conversely, if \(u_s\) is locally quasi-finite, then
\[
  u_X^{-1}\bigl(u_X(x)\bigr)
    \subset u_s^{-1}\bigl(u_s(x)\bigr)
\]
is a finite set. Since
\[
  u_X:X\longrightarrow H
\]
is locally of finite type, Chevalley's semicontinuity theorem
(EGA~IV\(_3\), 13.1.3) shows that there is an open neighborhood
\(W\) of \(x\) in \(X\) such that the fiber
\[
  u_X^{-1}\bigl(u_X(x')\bigr)
\]
is discrete for every \(x'\in W\). Hence \(u_X\) is quasi-finite
at \(x\), by EGA~II, 6.2.2 (and EGA~III\(_2\),
Err\(_{\mathrm{III}}\)~20).

Consequently,
\[
  U_{\mathrm{l.q.f.}}
    =\epsilon_G^{-1}\bigl(V_{\mathrm{l.q.f.}}\bigr),
  \qquad
  V_{\mathrm{l.q.f.}}
    =\pi_X^{-1}\bigl(U_{\mathrm{l.q.f.}}\bigr),
\]
and the first equality shows that \(U_{\mathrm{l.q.f.}}\) is open
in \(S\).

If, moreover, \(u\) is locally of finite type over \(S\), then
\(X=G\), and hence \(V_{\mathrm{l.q.f.}}\) is the inverse image
under \(\pi_G\) of \(U_{\mathrm{l.q.f.}}\).%
\begingroup
\renewcommand{\thefootnote}{\fnsymbol{footnote}}%
\footnote[1]{\label{note:sga3-VIB-2.5-u-loci}Translator's note:
in this sentence and in the analogous sentence for
\(V_{\mathrm{n.r.}}\) below, the French source prints \(G\) in
place of \(u\). The statement assumes that \(u\) is locally of
finite type (respectively, locally of finite presentation), and
\(X\) and \(Z\) are precisely the corresponding loci for \(u\);
the required \(u\) is used here.}%
\endgroup{}
This proves the first half of \textup{\SGARefLink{concl:sga3-VIB-2.5-ii}{(ii)}}.

Now consider the open subsets
\[
  V_{\mathrm{l.p.f.}}\supset V_{\mathrm{n.r.}},
\]
consisting of the points at which \(u\) is of finite presentation
(respectively, unramified), and suppose that
\[
  Z=V_{\mathrm{l.p.f.}}
\]
contains the identity section of \(G\). Let \(z\in Z\); set
\[
  s=\pi_Z(z),\qquad h=u_Z(z).
\]

If \(u\) is unramified at \(z\), then, by base change, \(u_s\) is
unramified at \(z\); and since \(G_s\) is assumed to be locally of
finite type, \(u_s\) is unramified by \SGARefLink{cor:sga3-VIB-1.3.1}{1.3.1}.

Conversely, if \(u_s\) is unramified at the point \(z\), then the
fiber
\[
  u_s^{-1}(h)=u^{-1}(h)
\]
is unramified over \(\kappa(h)\) at the point \(z\). Since \(Z\)
is an open subset of \(G\), the fiber
\[
  u_Z^{-1}(h)=u_s^{-1}(h)\cap Z
\]
is unramified over \(\kappa(h)\) at the point \(z\). Since \(u_Z\)
is locally of finite presentation, EGA~IV\(_4\), 17.4.1, implies
that \(u_Z\) is unramified at the point \(z\).

Thus
\[
  U_{\mathrm{n.r.}}
    =\epsilon_G^{-1}\bigl(V_{\mathrm{n.r.}}\bigr),
  \qquad
  V_{\mathrm{n.r.}}
    =\pi_Z^{-1}\bigl(U_{\mathrm{n.r.}}\bigr),
\]
and the first equality shows that \(U_{\mathrm{n.r.}}\) is open in
\(S\).

If, moreover, \(u\) is locally of finite presentation over \(S\),
then \(Z=G\), and hence \(V_{\mathrm{n.r.}}\) is the inverse image
under \(\pi_G\) of \(U_{\mathrm{n.r.}}\). This completes the proof
of assertion \textup{\SGARefLink{concl:sga3-VIB-2.5-ii}{(ii)}}.

\medskip
\phantomsection\label{cor:sga3-VIB-2.6}%
\noindent\textbf{Corollary 2.6.---}
\emph{Let \(u:G\to H\) be a morphism of \(S\)-groups that is
radicial (as is the case if \(u\) is a monomorphism;
EGA~I, 3.5.4). Suppose that \(G\) (respectively, \(H\)) is locally
of finite presentation and flat (respectively, locally of finite
type) over \(S\). Then the set \(U\) of \(s\in S\) such that
\(u_s\) is an open immersion is open in \(S\), and the restriction
of \(u\) over \(U\) is an open immersion.}

By \hyperref[concl:sga3-VIB-2.5-i]{2.5}\textup{\SGARefLink{concl:sga3-VIB-2.5-i}{(i)}}, the set \(U'\) of points \(s\in S\) such that
\(u_s\) is \'etale is open in \(S\). Since \(u\) is radicial, so
is \(u_s\), for every \(s\in S\); hence EGA~IV\(_4\), 17.9.1,
gives
\[
  U=U',
\]
which shows that \(U\) is open. Finally, by \hyperref[concl:sga3-VIB-2.5-i]{2.5}\textup{\SGARefLink{concl:sga3-VIB-2.5-i}{(i)}}, the
restriction of \(u\) over \(U\) is \'etale; since \(u\) is
radicial, this restriction is an open immersion
(EGA~IV\(_4\), 17.9.1).

\medskip
\phantomsection\label{prop:sga3-VIB-2.7}%
\noindent\textbf{Proposition 2.7.---}
\emph{Let \(G\) be an \(S\)-group. The following conditions are
equivalent:}
\begin{enumerate}
\item[\textup{\SGAInlineRefTarget{prop:sga3-VIB-2.7:item:i}(i)}] \emph{\(G\) is unramified over \(S\) at the
  points of the identity section;}
\item[\textup{\SGAInlineRefTarget{prop:sga3-VIB-2.7:item:ii}(ii)}] \emph{the identity section is an open
  immersion;}
\item[\textup{\SGAInlineRefTarget{prop:sga3-VIB-2.7:item:iii}(iii)}] \emph{\(G\) is of finite presentation over
  \(S\) at the points of the identity section, and, for every
  \(s\in S\), \(G_s\) is unramified over \(\kappa(s)\).}
\end{enumerate}
\emph{If, moreover, \(G\) is locally of finite presentation over
\(S\), then each of the three preceding conditions is equivalent
to the following:}
\begin{enumerate}
\item[\textup{(iv)}] \emph{\(G\) is unramified over \(S\).}
\end{enumerate}

The equivalence of assertions \textup{\SGARefLink{prop:sga3-VIB-2.7:item:i}{(i)}} and \textup{\SGARefLink{prop:sga3-VIB-2.7:item:ii}{(ii)}}
follows from the more general \hyperref[lem:sga3-VIB-2.7.1]{Lemma~2.7.1} below. Note (\SGARefLink{lem:sga3-VIB-1.3.1.1}{1.3.1.1})
that either condition \textup{\SGARefLink{prop:sga3-VIB-2.7:item:i}{(i)}} or condition \textup{\SGARefLink{prop:sga3-VIB-2.7:item:iii}{(iii)}}
implies that, for every \(s\in S\), \(G_s\) is locally of finite
type over \(\kappa(s)\). Thus, by EGA~IV\(_4\), 17.4.1, assertion
\textup{(i)} is equivalent to the assertion that, for every
\(s\in S\), \(G_s\) is unramified over \(\kappa(s)\) at the point
\(e_s\), the identity of \(G_s\), or, equivalently by \hyperref[cor:sga3-VIB-1.3.1]{1.3.1}, to
the assertion that \(G_s\) is unramified over \(\kappa(s)\).
Consequently, assertions \textup{\SGARefLink{cor:sga3-VIB-1.3.1:item:i}{(i)}} and \textup{\SGARefLink{prop:sga3-VIB-2.7:item:iii}{(iii)}} are
equivalent. Finally, if \(G\) is locally of finite presentation
over \(S\), assertions \textup{(iii)} and \textup{(iv)} are
equivalent (EGA~IV\(_4\), 17.4.1).

\medskip
\phantomsection\label{lem:sga3-VIB-2.7.1}%
\noindent\textbf{Lemma 2.7.1.---}
\emph{Let \(G\) be an \(S\)-scheme equipped with a section
\(\epsilon\). In order that \(G\) be unramified over \(S\) at the
points of this section, it is necessary and sufficient that
\(\epsilon\) be an open immersion.}

The condition is necessary by EGA~IV\(_4\), 17.4.1,
\textup{a)}\(\Rightarrow\)\textup{b\(^{\prime\prime}\)}).
Conversely, if \(\epsilon\) is an open immersion, then the
restriction to \(\epsilon(S)\) of the structural morphism
\[
  G\longrightarrow S
\]
is an isomorphism; hence \(G\) is unramified over \(S\) at the
points of \(\epsilon(S)\).
